\chaptertitle{第一章\quad 数据的收集与整理}

你每天早上走进教室，看到黑板上的值日表——"今天谁擦黑板？今天谁扫地？"值日表就是一种数据整理的结果：它把全班同学每天的劳动分配记录下来，让人一目了然。在生活中，我们经常需要收集数据、整理数据，然后从中发现规律或做出判断。这一章就来看看数据是怎么被收集和整理的。

\sectiontitle{1.1\quad 普查与抽样调查}

假如你是班长，想知道全班同学最喜欢哪种水果，好在下次班会上买给大家吃。最直接的办法是什么？一个一个问——问遍全班每一个人。这种对\textbf{所有}对象逐一进行调查的方式，叫做\textbf{全面调查}，也叫\textbf{普查}。

\begin{example}
判断以下情况中，普查是否可行：
(1) 检查一批灯泡的使用寿命；
(2) 了解全校$2000$名学生的近视情况；
(3) 检验一箱鸡蛋里有没有坏蛋。
\end{example}
\begin{solution}
(1) 不行。要测寿命就得把灯泡一直亮到坏掉——测完了灯泡也报废了。这种\textbf{破坏性检测}不能用普查。\\
(2) 可行，但要查遍$2000$人，工作量很大。如果只需要大致了解，可以只查一部分样本。\\
(3) 可行，但打开所有鸡蛋会造成浪费。一般做法是抽样检查几个。
\end{solution}

什么时候不能用普查？再看几个例子。

\begin{example}
以下场景应该用普查还是抽样调查？\\
(1) 人口普查；\\
(2) 某品牌方便面的口味满意度调查；\\
(3) 医院给新生婴儿做身体检查；\\
(4) 质检员测试一批钢缆的最大承受拉力。
\end{example}
\begin{solution}
(1) 普查。人口普查要求不漏一人。\\
(2) 抽样调查。不可能让全国人民都试吃。\\
(3) 普查。每个新生儿都必须检查。\\
(4) 抽样调查。拉断了的钢缆就没法卖了——又是破坏性检测。
\end{solution}

把这些规律总结一下：

\knowbox{
\textbf{普查}：对全体对象逐一调查。优点：结果准确、全面。缺点：费时费力，有时不可行。\\
\textbf{抽样调查}：从总体中抽取一部分进行调查，再推断整体。优点：省时省力。缺点：结果有误差。\\
\textbf{总体}：所要调查对象的全体。\quad \textbf{个体}：总体中的每一个对象。\\
\textbf{样本}：从总体中抽取出来的一部分个体。\\
\textbf{何时用抽样}：对象太多、调查具有破坏性、只需要大致了解情况时。}

\begin{center}
\begin{tikzpicture}[scale=0.9]
\draw[mainblue,thick] (0,0) ellipse (3 and 1.8);
\node at (0,1.2){总体（全校2000人）};
\filldraw[lightblue] (0,0) circle (1.2 and 0.8);
\draw[mainblue,thick] (0,0) circle (1.2 and 0.8);
\node at (0,0){样本（100人）};
\draw[->,thick,exred] (3.5,0.5) -- (1.5,0.5);
\node[exred] at (4.8,0.5){用样本推断总体};
\node at (0,-2.2){\small 图1-1\quad 总体与样本的关系};
\end{tikzpicture}
\end{center}

\begin{exercise}
\item 下列调查适合用普查还是抽样调查？说明理由。\\
(a) 检查一架飞机所有零件的合格情况；\\
(b) 调查某城市居民的平均月收入；\\
(c) 检测一批疫苗的有效期；\\
(d) 检查一批火柴能否划燃。
\item 某工厂生产了$10000$个螺丝，质检员从中抽取$200$个进行检测。这里的总体、个体、样本分别是什么？为什么要用抽样？
\item 有人说"抽样调查肯定不准，普查才准"，你同意吗？举一个反例说明你的观点。
\end{exercise}

\sectiontitle{1.2\quad 数据的整理——频数与频率}

收集到数据之后，一大堆数字摆在那里，看不出什么名堂。我们需要把它们整理成表格。

来看一个具体的例子。小明调查了全班$40$名同学最喜欢的运动项目，得到下面这些回答：

\begin{center}
足球，篮球，乒乓球，足球，羽毛球，篮球，足球，乒乓球，\\
篮球，足球，羽毛球，足球，乒乓球，篮球，羽毛球，足球，\\
篮球，足球，足球，乒乓球，篮球，羽毛球，足球，篮球，\\
乒乓球，足球，篮球，羽毛球，足球，乒乓球，足球，篮球，\\
羽毛球，足球，篮球，乒乓球，足球，羽毛球，篮球，足球
\end{center}

$40$个名字乱糟糟的，根本看不出哪项最受欢迎。如果用"正"字来统计每个项目出现的次数：

\begin{center}
\begin{tikzpicture}
\node at (0,0){
\begin{tabular}{c|c|c|c}
运动项目 & 画"正"字 & 频数（人数） & 频率\\\hline
足球 & 正正正正正$\,$— & 16 & $\dfrac{16}{40}=0.4$\\
篮球 & 正正正$\,$— & 11 & $\dfrac{11}{40}=0.275$\\
乒乓球 & 正$\,$正$\,$— & 7 & $\dfrac{7}{40}=0.175$\\
羽毛球 & 正$\,$— & 6 & $\dfrac{6}{40}=0.15$\\\hline
合计 & & 40 & 1
\end{tabular}
};
\node at (0,-2.5){\small 表1-1\quad 最喜欢的运动项目统计表};
\end{tikzpicture}
\end{center}

从这张表一眼就能看出：足球最受欢迎（16人），羽毛球最少（6人）。\textbf{频数}就是每个类别出现的次数，\textbf{频率}就是频数除以总数。注意各组频率加起来正好等于$1$。

\knowbox{
\textbf{频数}：每个类别出现的次数。\\
\textbf{频率}：频数$\div$总数。\\
性质：(1) 各组的频数之和等于总数；\\
\qquad (2) 各组的频率之和等于$1$。}

\begin{example}
某班$40$人期末数学成绩分五档：90分以上8人，80~89分15人，70~79分10人，60~69分5人，60分以下2人。求各档的频率，并验证频率之和为$1$。
\end{example}
\begin{solution}
总人数$=8+15+10+5+2=40$。\\
90分以上：$\dfrac{8}{40}=0.2$；\quad 80~89：$\dfrac{15}{40}=0.375$；\quad 70~79：$\dfrac{10}{40}=0.25$；\\
60~69：$\dfrac{5}{40}=0.125$；\quad 60以下：$\dfrac{2}{40}=0.05$。\\
频率之和$=0.2+0.375+0.25+0.125+0.05=1$。✓
\end{solution}

\begin{example}
一个样本$50$个数据分成$5$组，前四组的频数分别为$8,12,10,14$，求第五组的频数和频率。
\end{example}
\begin{solution}
第五组频数$=50-(8+12+10+14)=6$，频率$=\dfrac{6}{50}=0.12$。
\end{solution}

\begin{exercise}
\item 某班$40$人，统计喜欢的颜色：红12人，蓝10人，绿8人，黄6人，其他4人。制作频数-频率表。
\item 根据上题，红和蓝的频数之和是多少？频率之和呢？
\item 某校调查$200$名学生上下学方式：步行80人，公交60人，家长接送40人，骑车20人。求每种方式的频率，并画一个简单的频数分布表。
\end{exercise}

\sectiontitle{1.3\quad 抽样方法与调查设计}

抽样调查听起来简单，但实际操作中有很多讲究。如果抽样的方法不对，得到的样本可能不代表总体——这叫\textbf{样本偏差}。

最经典的例子：1936年美国总统大选，一家杂志寄出了$1000$万份问卷，回收$240$万份，预测共和党 candidate 会大胜。结果民主党获胜。为什么？因为问卷是从电话簿和汽车登记册上找地址的——1936年能装电话、有汽车的人偏富裕，倾向于投共和党。样本虽然大，但\textbf{有偏差}。

\begin{example}
以下抽样方法是否合理？为什么？\\
(1) 要调查全校学生的平均身高，只在篮球队里测量。\\
(2) 要了解市民对公交服务的满意度，只在早高峰的公交车上调查。\\
(3) 从全校$1000$人的学号中，随机抽出$50$个学号进行调查。
\end{example}
\begin{solution}
(1) 不合理。篮球队的同学普遍偏高，样本有偏差。\\
(2) 不合理。早高峰坐公交的人本来就对公交有依赖，而且心情可能烦躁，不能代表全体市民。\\
(3) 合理。这是\textbf{简单随机抽样}——每个学生被抽到的机会相等，样本有代表性。
\end{solution}

几种常用的抽样方法：

\knowbox{
\textbf{简单随机抽样}：每个个体被抽到的概率相等（如抽签、随机数表）。\\
\textbf{分层抽样}：将总体按某种特征分成若干层，在各层中按比例抽取。\\
\textbf{系统抽样}：将总体编号，按固定间隔抽取（如每$10$人抽$1$人）。\\
\textbf{便利抽样}：就近选取容易接触到的对象（结果可能有偏差，应尽量避免）。}

\begin{example}
某学校有男生$600$人、女生$400$人。要调查学生对校服款式的满意度，从全校$1000$人中抽$100$人。下面两种方案哪个更好？\\
方案A：随机抽$100$人。方案B：从男生中抽$60$人、女生中抽$40$人。
\end{example}
\begin{solution}
方案B更好。它保持了男女比例与总体一致（$6:4$），属于\textbf{分层抽样}。方案A是简单随机抽样，理论上可以，但如果恰好多抽了男生或少抽了女生，结果可能有偏差。分层抽样能保证各层比例准确。
\end{solution}

除了抽样方法，问卷的设计也大有学问。同样的问题，换一种问法，可能得到完全不同的答案。

\begin{example}
以下两组问题，你认为哪一组得到的答案更可靠？\\
(a) "你是否支持政府增加环保投入？"\\
(b) "你是否支持政府为了增加环保投入而适当提高税收？"\\
你觉得这两题的答案会一样吗？
\end{example}
\begin{solution}
很可能不一样。(a) 题大多数人会说"支持"，因为环保是好事。(b) 题提到了"提高税收"——一涉及自己掏钱，支持率就会下降。问卷的措辞会影响答案，设计问卷时要尽量\textbf{中立、无诱导性}。
\end{solution}

\begin{exercise}
\item 以下抽样方法合理吗？说明理由。\\
(a) 在电影院门口调查"市民喜欢看什么类型的电影"；\\
(b) 按班级学号，抽学号末位是$5$的同学调查作业量；\\
(c) 调查某河流的水质，只在排污口下游取样。
\item 你想了解"本校同学每天平均使用手机的时间"，设计一个抽样方案。
\item 为什么问卷中不应使用"很多人认为\ldots\ldots你是否也认为\ldots\ldots"这样的措辞？
\end{exercise}

\sectiontitle{本章小结}

本章我们学了数据从无到有的过程。首先决定怎么收集数据：普查（问遍所有人）还是抽样调查（只问一部分）。当对象太多、或者检查会损坏对象时，应该用抽样调查。拿到原始数据后，用频数和频率来整理——频数是"出现了几次"，频率是"频数除以总数"。抽样时要注意样本的代表性，避免样本偏差。问卷设计也要保持中立，不能诱导回答。从下一章开始，我们要把这些整理好的数据画成漂亮的图表。

\sectiontitle{课后练习}

\subsection*{A组\quad 基础巩固}

\begin{exercise}
\item 下列调查适合用普查还是抽样调查？\\
(a) 调查全国中小学生的平均身高；\\
(b) 检查运载火箭上每一个零件的质量；\\
(c) 测试某种新药的疗效；\\
(d) 检查超市里一批牛奶的保质期是否达标。
\item 某校有$1500$名学生，随机抽取$100$名调查每周阅读时间。指出总体、个体、样本。
\item 一个样本$60$个数据，各组的频数分别为$10,15,20,m$，求$m$和对应组的频率。
\item $40$名学生数学成绩：优秀（$\ge90$）$6$人，良好（$80$~$89$）$14$人，合格（$60$~$79$）$16$人，不合格（$<60$）$4$人。计算各等级的频率。
\item 已知某组数据的频率是$0.25$，频数是$12$，那么数据总数是多少？
\item 某班$50$人中，喜欢阅读的有$18$人。求喜欢阅读的频率和对应的扇形圆心角度数。
\end{exercise}

\subsection*{B组\quad 挑战提升}

\begin{exercise}
\item 某地区有$50000$户居民，要调查平均每户年用水量。设计一个合理的抽样方案（包括抽样方法和样本容量），并说明为什么你的方案比普查更合理。
\item 一个口袋中有红、黄、蓝三种颜色的球共$50$个。小明随机摸了$20$次（每次摸完放回），结果红球$8$次、黄球$7$次、蓝球$5$次。根据这些频率，估计三种球各有多少个？
\item 某校要调查学生睡眠时间。方案一：让班主任在班上问"昨天你睡了几个小时"。方案二：让学生匿名填写问卷。哪个方案更可靠？为什么？
\item 1936年美国大选投票前，某杂志寄出$1000$万份问卷，回收$240$万份，预测共和党获胜，结果民主党赢了。除了电话簿和汽车登记册的偏差外，还有一个问题：回收率只有$24\%$——回答问卷的人本身就不是随机样本。请解释为什么低回收率会导致偏差。
\item 要调查"中学生是否应该带手机上学"，设计两个版本的问题——一个有诱导性，一个中立——并说明为什么中立版本更可靠。
\end{exercise}

\newpage
